% adm-lapse-shift.tex — Lapse, shift, and time vector decomposition
% Inputable TikZ fragment for fig:adm-lapse-shift
% Requires: tikz with calc, arrows.meta (loaded by ehbook.cls)
%
% Schematic spacetime diagram showing two neighbouring spatial slices
% Sigma_t and Sigma_{t+dt}.  From point P on the lower slice, three
% vectors illustrate the ADM decomposition:
%   - alpha n^mu (normal, vertical)   — EHJetBlue
%   - (d/dt)^mu  (time vector, tilted) — EHHorizonCrimson
%   - beta^i dt  (shift, tangential)   — EHAccretionGold
% The shift connects where the normal arrives (Q) to where the time
% vector arrives (R) on the upper slice.
%
% Sheet coordinates reuse the foliation diagram's perspective layout:
%   Front-to-back offset: dx = -0.7, dy = +1.5
%   Lower: FL=(0.5,0) FR=(7.0,0.20) BR=(6.3,1.70) BL=(-0.2,1.50)
%   Upper: FL=(0.5,2.50) FR=(7.0,3.40) BR=(6.3,4.90) BL=(-0.2,4.00)

\begin{tikzpicture}[
    >={Stealth[length=5pt]},
    font=\footnotesize,
  ]

  % ── Lower sheet (Σ_t) ─────────────────────────────────────────────

  \fill[EHMarginGray, opacity=0.06]
    plot[smooth, tension=0.6] coordinates {
      (0.5, 0.0) (2.5, -0.08) (4.5, 0.06) (7.0, 0.20)}
    -- (6.3, 1.70) -- (-0.2, 1.50) -- cycle;

  \draw[EHMarginGray, semithick]
    plot[smooth, tension=0.6] coordinates {
      (0.5, 0.0) (2.5, -0.08) (4.5, 0.06) (7.0, 0.20)};

  \draw[EHMarginGray, thin, opacity=0.4]
    (-0.2, 1.50) -- (6.3, 1.70);
  \draw[EHMarginGray, thin, opacity=0.4]
    (0.5, 0.0) -- (-0.2, 1.50)  (7.0, 0.20) -- (6.3, 1.70);

  % Coordinate gridlines — transverse
  \draw[EHMarginGray, very thin, opacity=0.15]
    (0.27, 0.50) -- (6.77, 0.70)
    (0.03, 1.00) -- (6.53, 1.20);

  % Coordinate gridlines — longitudinal
  \draw[EHMarginGray, very thin, opacity=0.15]
    (2.13, 0.05) -- (1.43, 1.55)
    (3.75, 0.10) -- (3.05, 1.60)
    (5.38, 0.15) -- (4.68, 1.65);

  \node[font=\small, text=EHMarginGray, anchor=west]
    at (7.15, 0.20) {$\Sigma_t$};


  % ── Upper sheet (Σ_{t+dt}) ────────────────────────────────────────

  \fill[EHMarginGray, opacity=0.06]
    plot[smooth, tension=0.6] coordinates {
      (0.5, 2.50) (2.5, 2.44) (4.5, 2.96) (7.0, 3.40)}
    -- (6.3, 4.90) -- (-0.2, 4.00) -- cycle;

  \draw[EHMarginGray, semithick]
    plot[smooth, tension=0.6] coordinates {
      (0.5, 2.50) (2.5, 2.44) (4.5, 2.96) (7.0, 3.40)};

  \draw[EHMarginGray, thin, opacity=0.4]
    (-0.2, 4.00) -- (6.3, 4.90);
  \draw[EHMarginGray, thin, opacity=0.4]
    (0.5, 2.50) -- (-0.2, 4.00)  (7.0, 3.40) -- (6.3, 4.90);

  % Gridlines — transverse
  \draw[EHMarginGray, very thin, opacity=0.15]
    (0.27, 3.00) -- (6.77, 3.90)
    (0.03, 3.50) -- (6.53, 4.40);

  % Gridlines — longitudinal
  \draw[EHMarginGray, very thin, opacity=0.15]
    (2.13, 2.73) -- (1.43, 4.23)
    (3.75, 2.95) -- (3.05, 4.45)
    (5.38, 3.18) -- (4.68, 4.68);

  \node[font=\small, text=EHMarginGray, anchor=west]
    at (7.15, 3.40) {$\Sigma_{t+\mathrm{d}t}$};


  % ── Point P on lower sheet ────────────────────────────────────────
  \coordinate (P) at (3.0, 0.55);

  % ── Points on upper sheet ─────────────────────────────────────────
  % Q: where the Eulerian observer arrives (along n^μ, directly above P)
  \coordinate (Q) at (3.0, 3.20);
  % R: where the coordinate time vector arrives (same labels x^i)
  \coordinate (R) at (4.8, 3.40);


  % ── Unit normal α n^μ (EHJetBlue) ────────────────────────────────
  \draw[->, EHJetBlue, very thick, shorten >=2pt]
    (P) -- (Q);
  \node[font=\scriptsize, text=EHJetBlue, anchor=east, inner sep=2pt]
    at (2.95, 1.90) {$\alpha\,n^\mu$};


  % ── Time vector (∂/∂t)^μ (EHHorizonCrimson) ──────────────────────
  \draw[->, EHHorizonCrimson, very thick, shorten >=2pt]
    (P) -- (R);
  \node[font=\scriptsize, text=EHHorizonCrimson, anchor=west, inner sep=2pt]
    at (4.20, 1.85) {$(\partial/\partial t)^\mu$};


  % ── Shift displacement β^i dt (EHAccretionGold) ──────────────────
  \draw[->, EHAccretionGold, very thick, shorten >=2pt]
    (Q) -- (R);
  \node[font=\scriptsize, text=EHAccretionGold, anchor=south, inner sep=2pt]
    at (3.9, 3.38) {$\beta^i\,\mathrm{d}t$};


  % ── Right-angle marker at Q ──────────────────────────────────────
  %
  % Shows orthogonality between n^μ (vertical) and β^μ (tangential).
  \draw[thin] (3.0, 3.05) -- (3.15, 3.07) -- (3.15, 3.20);


  % ── Point markers (drawn last, on top) ───────────────────────────
  \fill (P) circle (1.5pt);
  \node[font=\scriptsize, anchor=north east, inner sep=2pt]
    at (P) {$P$};

  \fill (Q) circle (1.5pt);
  \fill (R) circle (1.5pt);

\end{tikzpicture}
